% frontmatter/source-map.tex

\addchap{Карта первоисточников}
\label{ch:primary-sources}
Когда в платёжной команде звучит фраза \enquote{надо проверить
в правилах}, спор часто начинается об источнике истины. Один человек открывает свод правил сети, другой ищет ответ
в технической спецификации, третий пересказывает заметку поставщика.
У одной и той же темы действительно могут существовать три слоя
документов: публичные правила для участников, технические
спецификации и закрытые оперативные уведомления.

Поэтому первая задача навигации проста: прежде чем читать, нужно
понять, какой документ в этой теме главный, а какой служит лишь
комментарием, пересказом или удобным порталом для интеграции. Эта глава помогает быстрее выбрать
правильную точку входа в первоисточники.

\begin{table}[h]
  \centering
  \small
  \setlength{\tabcolsep}{4pt}
  \renewcommand{\arraystretch}{1.15}
  \begin{tabularx}{\textwidth}{
    >{\RaggedRight\arraybackslash}p{2.0cm}
    >{\RaggedRight\arraybackslash}X
    >{\RaggedRight\arraybackslash}p{2.8cm}
    >{\RaggedRight\arraybackslash}X
  }
    \textbf{Организация} & \textbf{Ключевые документы}
    & \textbf{Доступ}     & \textbf{Когда читать}
    \\
    EMVCo                & Books 1--4, Contactless, Tokenisation, 3DS
    & в основном открытый & когда нужно понять логику карты,
    терминала, токена или 3DS      \\
    PCI SSC              & PCI DSS, SAQ, PIN Security, P2PE
    & открытый            & когда вопрос о границах PCI-контура,
    контролях или аудите       \\
    Visa                 & Core Rules, dispute guides, TADG, BASE
    I/II, VTS      & смешанный           & когда вопрос об
    обязанностях участников или потоке VisaNet      \\
    Mastercard           & Mastercard Rules, Chargeback Guide,
    TPR/Banknet, MDES & смешанный           & когда нужны Banknet,
    IPM или правила Mastercard                 \\
    НСПК / Мир           & Правила Мир, документы СБП,
    техспецификации НСПК      & в основном закрытый & когда вопрос о
    российском контуре, Мир или СБП                  \\
    Банк России          & 161-ФЗ, 762-П, 821-П, 851-П, 266-П
    & открытый            & когда вопрос про правовую рамку и
    обязательные требования       \\
    ISO                  & ISO 8583, ISO 20022, ISO 4217, ISO 9564
    & смешанный / платный & когда нужен язык стандарта и общая
    модель, а не сетевой диалект \\
  \end{tabularx}
\end{table}

Каждый корпус документов разобран ниже: с какого документа начинать
и как не перепутать правило участника с протоколом
или вспомогательным порталом.

\section{EMVCo}
\label{sec:ps-emvco}

EMVCo нужен там, где спор идёт не о договорных обязанностях,
а о поведении самой платёжной техники. Если вопрос звучит как
\enquote{что именно должны сделать терминал, карта, токен
или мобильный кошелёк}, искать ответ нужно прежде всего здесь.
Ни правила эквайера, ни правила карточной сети этот слой
не заменят.

\subsection*{Какие документы являются основными}

Если речь идёт о контактной EMV-транзакции, точкой входа остаются
\textbf{EMV Chip Specifications}, Books 1--4
(см. главу \ref{ch:emv}). Для бесконтактной оплаты нужен корпус
\textbf{EMV Contactless Specifications}, где описаны NFC
и логика бесконтактного ядра
(см. главу \ref{ch:contactless}). Вопросы токенизации закрывает
\textbf{EMV Payment Tokenisation Specification}: роли TSP,
ограничения домена токена и его жизненный цикл
(см. главу \ref{ch:tokenization}). Для 3DS основным документом
служит \textbf{EMV 3-D Secure Protocol and Core Functions
Specification}, где описаны сообщения AReq/ARes, сценарий
дополнительной проверки и общая логика протокола
(см. главу \ref{ch:3ds}).
\begin{itemize}
  \item \fullcite{emvco-specs}
  \item \fullcite{emvco-contactless-chip}
  \item \fullcite{emvco-payment-tokenisation}
  \item \fullcite{emvco-3ds}
\end{itemize}

\subsection*{Как входить}

EMVCo хуже всего читается подряд. Здесь почти всегда выгоднее идти
от вопроса к разделу. Если нужно понять жизненный цикл контактной
операции, полезнее всего Book 3. Если важно поведение терминала,
кассира и интерфейса приёма, чаще открывают Book 4. Если вопрос
относится к бесконтактному входу в транзакцию, нужен соответствующий
том Contactless. Иными словами, EMVCo редко читают от корки
до корки, зато именно эти документы лучше всего подходят
для точечной проверки спорного места.

\section{PCI SSC}
\label{sec:ps-pci}

PCI SSC нужен тогда, когда команда спорит о границах PCI-контура
и наборе обязательных контролей. Можно ли ограничиться SAQ A, должен ли
мерчант контролировать JavaScript на странице оплаты, как защищается
PIN block и где проходит точка дешифрования, -- все эти вопросы
нужно закрывать документами PCI, а не пересказами интеграторов.

\subsection*{Какие документы являются основными}

\textbf{PCI DSS} остаётся главным документом по контролям среды,
где появляются данные держателя карты
(см. главу \ref{ch:pci-dss}). \textbf{SAQ} нужны как сокращённые
режимы оценки для разных моделей интеграции.
\textbf{PCI PIN Security Requirements} отвечают за работу
с PIN и HSM. \textbf{PCI P2PE} нужен тогда, когда вопрос
упирается в шифрование от терминала до точки дешифрования.
\begin{itemize}
  \item \fullcite{pci-dss-4}
  \item \fullcite{pci-saq}
  \item \fullcite{pci-pin-security}
  \item \fullcite{pci-p2pe}
\end{itemize}

\subsection*{Как входить}

Практический порядок чтения почти всегда один и тот же: сначала
определить границы PCI-контура, потом понять, какой путь оценки
применим, и лишь затем разбирать отдельные требования. Ошибка
в определении границ среды обычно обходится дороже, чем неверное
чтение одного подпункта.

\section{Visa}
\label{sec:ps-visa}

У Visa особенно легко перепутать два слоя: правила
для участников и техническую документацию VisaNet. Первый слой
отвечает на вопрос \enquote{что разрешено и кто обязан}, второй --
на вопрос \enquote{как это устроено в сообщениях, терминалах
и сервисах сети}.

\subsection*{Какие документы являются основными}

\textbf{Visa Core Rules and Visa Product and Service Rules}
нужны, когда речь идёт об обязанностях эмитента, эквайера
и мерчанта. \textbf{Dispute Management Guidelines for Visa Merchants}
-- удобная публичная точка входа в диспутную логику; если
возникает расхождение, приоритет всё равно остаётся за Core Rules.
\textbf{Visa Technical Acceptance Device Guide (TADG)}
описывает приём карт на терминалах. Для авторизации и клиринга
в VisaNet ориентиром служат материалы \textbf{BASE I / BASE II};
публичные ориентиры по ним есть у Visa, но рабочая техническая
детализация обычно доступна участникам сети.
\textbf{VTS documentation} закрывает вопросы Visa Token Service.
\begin{itemize}
  \item \fullcite{visa-core-rules}
  \item \fullcite{visa-dispute-guidelines-merchants}
  \item \fullcite{visa-tadg}
  \item \fullcite{visa-visanet-booklet}
  \item \fullcite{visa-vts}
\end{itemize}

\subsection*{В каком порядке читать}

Если спор касается ролей участников или допустимости сценария,
начинать нужно с Core Rules. Если вопрос в поведении терминала
или в техническом поле на приёме, нужен TADG или закрытая
техническая спецификация через участника сети. Портал
для разработчиков \url{https://developer.visa.com} полезен
для API-продуктов; по сетевым правилам следует обращаться к Core Rules напрямую.

\section{Mastercard}
\label{sec:ps-mastercard}

У Mastercard композиция похожа, но названия документов
другие: отдельно живут правила и диспутная политика, отдельно --
технические материалы Banknet/IPM и сервисная документация
для разработчиков.

\subsection*{Какие документы являются основными}

\textbf{Mastercard Rules} задают обязанности участников.
\textbf{Mastercard Chargeback Guide} нужен для логики диспутов
и chargeback. \textbf{Transaction Processing Rules}
и связанные технические материалы Banknet/IPM служат точкой
входа в авторизацию и клиринг через Banknet.
\textbf{MDES} описывает токенизацию Mastercard.
\textbf{M/Chip Requirements} нужны, когда важно понять
Mastercard-специфику EMV.
\begin{itemize}
  \item \fullcite{mc-rules}
  \item \fullcite{mc-chargeback-guide}
  \item \fullcite{mc-tpr}
  \item \fullcite{mc-switching-services}
  \item \fullcite{mc-mdes-token-connect}
  \item \fullcite{mc-mchip}
\end{itemize}

\subsection*{В каком порядке читать}

Если вопрос договорный или операционный, обычно начинают
с Rules. Если нужен протокол или клиринговая модель, переходят
к материалам Banknet и IPM. Портал для разработчиков Mastercard
полезен для открытых API; закрытые технические документы
следует получать через участника сети напрямую.

\section{НСПК и Мир}
\label{sec:ps-nspk}

В российском контуре главная трудность в другом: рамка ролей
и продуктов частично видна в публичных правилах, а значительная
часть технических деталей скрыта в документах для участников.
Поэтому здесь особенно важно заранее понять, ищете вы правовую
рамку, продуктовую модель или конкретный формат сообщения.

\subsection*{Какие документы являются основными}

\textbf{Правила платёжной системы Мир} задают роли и обязанности
участников. \textbf{Спецификация бесконтактных платежей Мир}
нужна для бесконтактного приёма. \textbf{Технические спецификации
НСПК} описывают форматы сообщений и подключение. Для СБП публичной
точкой входа служат сайт СБП и документы Банка России по подключению;
подробные технические требования и тестовые материалы НСПК выдаёт
через портал поддержки участникам и банкам-кандидатам.
\begin{itemize}
  \item \fullcite{nspk-mir-rules}
  \item \fullcite{nspk-mir-contactless}
  \item \fullcite{nspk-tsp}
  \item \fullcite{sbp-main}
  \item \fullcite{cbr-sbp-actions}
\end{itemize}

\subsection*{Как входить}

Если вопрос касается роли участника, продуктовой схемы или общего
режима работы, обычно хватает публичных правил. Если же спор
упирается в формат сообщения, сертификацию или детали
внутрироссийского процессинга, почти наверняка понадобится
закрытая документация через участника системы или портал поддержки
НСПК.

\section{Банк России}
\label{sec:ps-cbr}

Банк России нужен тогда, когда требуется понять границы допустимого поведения.
Как только вопрос звучит как \enquote{обязаны ли мы},
\enquote{имеем ли мы право} или \enquote{какой статус нужен},
переходить следует к нормативным актам ЦБ и федеральным законам.

\subsection*{Какие документы являются основными}

161-ФЗ задаёт базовые роли платёжного рынка.
762-П описывает правила перевода денежных средств.
821-П и 851-П фиксируют текущие требования по защите
информации. 266-П сохраняет карточную специфику.
Для цифрового рубля, СБП и других специальных тем набор документов
будет своим, но логика поиска не меняется: начинать нужно
с нормативного текста, а пересказ использовать только
как вспомогательную опору.
\begin{itemize}
  \item \fullcite{fz-161}
  \item \fullcite{cbr-762p}
  \item \fullcite{cbr-821p}
  \item \fullcite{cbr-851p}
  \item \fullcite{cbr-266p}
\end{itemize}

\subsection*{Как входить}

Нормативные акты полезно читать связкой: закон задаёт роли
и пределы модели, положение описывает операционную логику,
письма и разъяснения показывают, как регулятор читает спорные
места. Для практической работы удобнее пользоваться правовыми
системами, где видны история изменений и перекрёстные ссылки.

\section{ISO}
\label{sec:ps-iso}

ISO нужен тогда, когда надо выйти на более общий язык, чем язык
конкретной сети. Это слой стандартов, на котором платёжная
индустрия договаривается о структуре сообщений и терминов, хотя
в реальной системе почти всегда поверх него живёт сетевой
диалект.

\subsection*{Какие документы являются основными}

ISO~8583 остаётся базовым стандартом по структуре финансовых
сообщений, хотя на практике его всегда конкретизируют сетевые
реализации. ISO~20022 задаёт современную модель финансовых
сообщений и особенно важен для трансграничных контуров, SWIFT MX
и многих систем мгновенных платежей
(см. главу \ref{ch:multicurrency}). ISO~4217, ISO~3166-1
и ISO~9564-1 служат базовыми документами по кодам валют,
кодам стран и форматам PIN block.
\begin{itemize}
  \item \fullcite{iso-8583-2023}
  \item \fullcite{iso-20022}
  \item \fullcite{iso-4217}
  \item \fullcite{iso-3166-1}
  \item \fullcite{iso-9564}
\end{itemize}

\subsection*{Как входить}

Если задача прикладная и привязана к конкретной сети, сначала
имеет смысл читать сетевой диалект, а потом возвращаться к ISO
как к верхнему слою терминов и модели. Если же задача
архитектурная, межсистемная или связана с проектированием
нового обмена, ISO~20022 полезно открывать сразу
как исходную модель сообщения.

\section*{Итоги}
\begin{itemize}
  \item Для каждой темы есть свой главный документ: EMVCo для
    поведения карты, терминала, токена и 3DS, PCI SSC
    для контролей, карточные сети для правил участников
    и собственных контуров, ЦБ для правовой рамки, ISO
    для языка стандартов.
  \item Открытые и закрытые документы стоит разводить заранее:
    это экономит время и не создаёт ложного ожидания,
    что любой вопрос можно закрыть одним публичным PDF.
  \item Начинать чтение лучше с документа, который отвечает
    именно на ваш тип вопроса: правила участников, протокол,
    контрольная рамка или нормативное требование.
  \item Порталы для разработчиков полезны для API-интеграции,
    но редко бывают главным источником по сетевым правилам.
  \item Хорошая навигация начинается с карты источников:
    кто в этой теме владеет истиной, где документ доступен
    и в каком порядке его читать.
\end{itemize}
